\begin{textsample}{NEG Dim 1 – human – Score -6.00 – mg/mg\_ef\_linguagens\_ed\_fis\_1\_p4.txt}  \label{ex:f1_neg_019}
A materialidade de nossa existência está sustentada pela nossa corporeidade : do momento da concepção até nosso último minuto de vida estamos vinculados aos \textbf{movimentos} cíclicos, acíclicos, voluntários, involuntários, conscientes, expressivos, etc. de partes do nosso corpo ou dele como um todo. \textbf{Produzimos} valores, saberes, conhecimentos, princípios, normas e conceitos a partir das expressões e relações de todas as nossas dimensões constitutivas.
Nossa sociedade ainda valoriza muito a dimensão mental/cognitiva como processo de aquisição e produção de conhecimento.
Mas diversas teorias \textbf{têm} oferecido legitimidade e importância para a participação de outras dimensões – motora/corporal, afetivo/emocional, social/cultural e espiritual/existenciais- neste processo.
Ou ainda mais, estas dimensões \textbf{têm} sido validadas como imprescindíveis para a formação humana e cidadã atualmente almejada, também expressa nas 10 Competências Gerais da Base Nacional Comum Curricular.
Respaldamos neste \textbf{sentido} as aulas de Educação Física na escola como tempo e espaço privilegiado que pode garantir o \textbf{reconhecimento} e a valorização da vivência destas dimensões e seu potencial formativo.

% matched lemmas: movimento, produzir, reconhecimento, sentido, ter
\end{textsample}
